\section{江面以下：孙吴的国家怎样抵达家庭}

\subsection{从建业到临湘，命令会变成什么}

建业的诏令在史书中往往以一两句结束，走马楼吴简却让命令抵达一座县廷以后
的样子变得繁复。木简上的户名、田地、租米、吏役和勾稽不是宏大制度的缩影，
而是临湘县某些机构留下的一次性工作痕迹：书手要抄录，保管者要收发，基层
中介要辨认户与田，拖欠或变更又要另作说明。国家并非先有一部完整税法才
来到乡里；它是在这些可以被写下、核对、追问的事务中逐渐获得重量。

正因为简牍是地方性残存，不能把它们换算成孙吴总人口或全境亩数。临湘的田租
形式不能自动代表建业、会稽或交州，保存下来的文书也偏向原本需要行政处理的
人和物。它的价值正在于阻止另一种误读：若只读《三国志·吴书》的帝王、战役
和册命，便容易以为江东的国家只由宫门和舰队组成。吴简告诉我们，水师的粮、
都城的役和地方官的责任，须落在许多不留姓名的笔画、仓门和往返之中。

户籍既能给国家以可征发的人，也能给家庭带来被辨认和被追责的风险。某户被
登记，并不等于其成员没有迁徙、逃役、依附或与邻里协商的余地；恰好相反，
反复核对的必要说明这种余地一直存在。没有完整的口数序列，我们不能为三国
江南制造精确人口曲线。可以确认的是，战争和江防所需的粮、人、船，不会只
在都城的军事会议里出现，它们必须经过县廷所能处理的分类和关系。

\subsection{钱币不是一个“经济背景”}

孙权铸大钱的记载看上去像一条简短的货币政策。把它单独视为“吴国经济”
就会失去它与军政压力的联系。铸钱需要金属、工匠、官府许可和人们愿意接受
它的交易场景；高面值是否真能便利支付，取决于市场、税收与物资供给，而不是
诏令写出的数字。《三国志》可让我们知道朝廷尝试过这项安排，却没有留下足以
重建兑换率、物价或各地使用范围的连续材料。

因此大钱应被读作一个正在试图调度资源的政府留下的痕迹，而不应被用来证明
吴“富”或“穷”。同样，宫室营建和役使争议不是可加总为徭役总数的账簿。
它们让我们看到中枢会把人力从田野、运输和江防中抽走；不同地方如何分担、
谁能避开、谁承担最重，现存材料没有足够均匀的记录。用不确定的总数填满
这种空白，会比承认空白更远离这些文书的尺度。

\subsection{船、港口与沿江的劳作}

孙吴被称为水军之国，容易让读者只看见江上的战船。船需要木材、铁器、绳索、
船工、划手、修理点和能卸下粮食的港口；一支在建业出发的队伍还要经过中游
军镇，才可能回应荆州或淮南的战事。卫温、诸葛直的远航之所以值得保留，不在
于为夷洲、亶洲找到现代地图上的唯一位置，而在于它暴露出水路能力的另一面：
离开熟悉的江岸后，风候、疾病、饮水和补给会迅速夺走朝廷自以为能够控制的
时间。

《三国志·吴主传》记录远航归来与处分，不能把船员死亡数和岛屿位置补成确数。
不过，失败本身改变了读法。孙吴并非因为拥有长江便天然拥有海洋；它的行动
要在守江、防魏、转运和探索之间竞争同样的材料与劳力。沿海群体也不能只作为
吴廷计划书里的“可得人口”出现。史书没有留下他们持续而独立的声音时，读者
应把沉默留在海图上，而不是以后来疆域替他们说话。

\subsection{文字、墓葬与不被写进本纪的人}

一块墓志、一座墓葬或一段宗教文本，所能照亮的通常是被纪念、被埋葬或被抄写
的人的世界，而非一座州郡的全部。它们仍十分重要：与官修帝纪相比，这些材料
让家族称谓、随葬习惯、信仰语言和地方书写偶尔显现出来。问题不是把它们降为
“补充材料”，而是问它们为什么会保存、为谁而作，以及沉默了谁。墓志中的
家世修辞不能直接证明实际财富，宗教叙事中的神异不能自动成为行政事实；但
两者都能提醒我们，三国人的死亡、记忆和救赎并不等候一次政权更替才发生。

石经残片与吴简的对照尤其有用。前者来自洛阳的官学与公共校读，后者来自江南
县廷的日常手续；一个让国家希望固定的经典文字可见，一个让国家实际处理的
人、田和米可见。它们不应被拼成“全国文化程度”的统计表，却共同冲淡了
魏本纪的单一视角。三国并不只由皇帝在首都做决定，也由文字怎样被刻、抄、
收发、保存和遗失来决定哪些生活后来仍能被我们看见。
